\chapter{IMPLEMENTATION}

\justify

This chapter describes the implementation methodology and key modules of the HPDC Defect Detection System. The implementation is divided into three functional stages: data engineering (patch generation, balancing, augmentation), model training (RF-DETR fine-tuning), and inference.

\section{Proposed Methodology}

The core methodology is a \textbf{patch-based sliding-window training pipeline} designed specifically for extreme small-object detection. Since crack bounding boxes occupy less than 0.1\% of the full image area, full-image training is infeasible. The approach divides each full HPDC image into overlapping 1024$\times$1024 patches, with Shapely polygon-aware annotation clipping to preserve accurate ground-truth boundaries.

\begin{figure}[H]
\centering
\begin{tikzpicture}[node distance=0.5cm,
block/.style={rectangle, rounded corners=4pt, draw=magnaDark,
fill=magnaLightGray, text width=7.5cm,
minimum height=1cm, align=center, font=\small},
arrow/.style={-Stealth, thick, color=magnaRed}]

\node[block, fill=magnaBlue!15, draw=magnaBlue] (input)
    {\textbf{Input:} Full image + COCO polygon annotations};
\node[block, below=0.5cm of input] (stride)
    {\textbf{Sliding Window:} Stride $= 1024 \times (1 - 0.35) = 665$\,px};
\node[block, below=0.5cm of stride] (extract)
    {\textbf{Patch Extraction:} Crop \texttt{image[y:y+H, x:x+W]}};
\node[block, below=0.5cm of extract] (shapely)
    {\textbf{Polygon Clipping (Shapely):} Intersect polygon with patch boundary;
    recompute bbox from clipped polygon};
\node[block, below=0.5cm of shapely] (balance)
    {\textbf{Balancing:} Downsample negatives to 1:2 pos:neg ratio};
\node[block, below=0.5cm of balance] (aug)
    {\textbf{Augmentation:} Horizontal flip on crack patches only};
\node[block, fill=successGreen!10, draw=successGreen, below=0.5cm of aug] (output)
    {\textbf{Output:} COCO-format JSON + PNG patches};

\draw[arrow] (input)   -- (stride);
\draw[arrow] (stride)  -- (extract);
\draw[arrow] (extract) -- (shapely);
\draw[arrow] (shapely) -- (balance);
\draw[arrow] (balance) -- (aug);
\draw[arrow] (aug)     -- (output);
\end{tikzpicture}
\caption{Proposed methodology flowchart: patch-based data engineering pipeline.}
\label{fig:methodology}
\end{figure}

The EMA (Exponential Moving Average) training strategy smooths weight updates across epochs:
\begin{equation}
\boldsymbol{\theta}_{\text{EMA}} = \alpha \cdot \boldsymbol{\theta}_{\text{EMA}} + (1-\alpha) \cdot \boldsymbol{\theta}_{\text{current}}
\end{equation}
where $\alpha$ is the decay factor (typically 0.9999). EMA weights consistently outperformed instantaneous weights across all experiments.

\section{Description of Modules}

\subsection*{Module 1: Patch Generator}

\begin{itemize}[leftmargin=2em, topsep=2pt]
    \item \textbf{Input:} Full HPDC image (4096$\times$3000 PNG), COCO annotation JSON
    \item \textbf{Output:} Set of 1024$\times$1024 PNG patches + updated COCO JSON with clipped annotations
    \item \textbf{Description:} Enumerates sliding-window positions with stride 665 px. For each patch, clips all polygon annotations using Shapely geometry intersection. Recomputes axis-aligned bounding boxes from clipped polygons. Discards annotations where intersection area falls below a minimum threshold.
\end{itemize}

\begin{lstlisting}[language=Python, caption={Patch Generator pseudocode}]
def generate_patches(image, annotations, patch_size=1024, overlap=0.35):
    stride = int(patch_size * (1 - overlap))   # = 665 px
    patches = []
    for y in range(0, image.height - patch_size + 1, stride):
        for x in range(0, image.width - patch_size + 1, stride):
            patch = image.crop(x, y, patch_size, patch_size)
            patch_box = Polygon([(x,y),(x+S,y),(x+S,y+S),(x,y+S)])
            clipped_anns = []
            for ann in annotations:
                poly = Polygon(ann['segmentation'])
                intersection = poly.intersection(patch_box)
                if intersection.area > MIN_AREA_THRESHOLD:
                    new_bbox = intersection.bounds  # recomputed bbox
                    clipped_anns.append(make_coco_ann(new_bbox, ann['category_id']))
            patches.append((patch, clipped_anns))
    return patches
\end{lstlisting}

\subsection*{Module 2: Dataset Balancer and Augmentor}

\begin{itemize}[leftmargin=2em, topsep=2pt]
    \item \textbf{Input:} List of all generated patches with their annotation counts
    \item \textbf{Output:} Balanced and augmented patch dataset
    \item \textbf{Description:} Separates patches into positive (contains $\geq 1$ defect annotation) and negative (background only). Randomly downsamples negatives to achieve a 1:2 positive:negative ratio. Then applies horizontal flip to crack-containing patches, recomputing bounding box coordinates as: $x_{\text{new}} = W - x - w$.
\end{itemize}

\begin{lstlisting}[language=Python, caption={Balancing and augmentation pseudocode}]
def balance_and_augment(patches):
    positives = [p for p in patches if len(p.annotations) > 0]
    negatives  = [p for p in patches if len(p.annotations) == 0]
    # Downsample negatives to 2x positives
    negatives = random.sample(negatives, min(len(negatives), 2*len(positives)))
    dataset = positives + negatives
    # Augment crack patches with horizontal flip
    augmented = []
    for patch in dataset:
        if has_crack(patch):
            flipped_img  = cv2.flip(patch.image, 1)
            flipped_anns = [flip_bbox(ann, patch.width) for ann in patch.annotations]
            augmented.append(Patch(flipped_img, flipped_anns))
    return dataset + augmented

def flip_bbox(ann, W):
    x, y, w, h = ann['bbox']
    return {'bbox': [W - x - w, y, w, h], 'category_id': ann['category_id']}
\end{lstlisting}

\subsection*{Module 3: RF-DETR Trainer}

\begin{itemize}[leftmargin=2em, topsep=2pt]
    \item \textbf{Input:} COCO-format patch dataset (train/val/test splits), training configuration
    \item \textbf{Output:} Best EMA checkpoint (\texttt{checkpoint\_best\_ema.pth}), per-epoch mAP metrics
    \item \textbf{Description:} Fine-tunes RF-DETR Medium with DINOv2 ViT-B backbone. Training uses Hungarian bipartite matching loss (classification + $\mathcal{L}_1$ + GIoU). EMA weights are maintained throughout and saved when validation EMA mAP50-95 improves.
\end{itemize}

\begin{lstlisting}[language=Python, caption={Training configuration (Experiment 7)}]
model = RFDETRMedium(pretrain_weights='rf-detr-medium.pth')
model.train(
    dataset_dir   = 'RF_DETR_Medium_1024/',
    epochs        = 20,
    batch_size    = 12,
    imgsz         = 1024,
    workers       = 16,
    device        = 'cuda'
)
# Best checkpoint: checkpoint_best_ema.pth
# Best EMA mAP50-95: 0.2379  |  mAP50: 0.4647
\end{lstlisting}

\subsection*{Module 4: Inference Engine}

\begin{itemize}[leftmargin=2em, topsep=2pt]
    \item \textbf{Input:} Full HPDC casting image, trained EMA checkpoint
    \item \textbf{Output:} Annotated image with bounding boxes for Crack / Cold Shut detections
    \item \textbf{Description:} Loads the EMA checkpoint, generates 1024$\times$1024 patches from the input image, runs RF-DETR forward pass on each patch, maps bounding box coordinates back to the full image coordinate system, and renders annotated output.
\end{itemize}

\begin{lstlisting}[language=Python, caption={Inference pseudocode}]
def detect(image_path, checkpoint, conf_threshold=0.3):
    model = RFDETRMedium()
    model.load_state_dict(torch.load(checkpoint))
    model.eval()
    image   = load_image(image_path)
    patches = generate_patches(image, patch_size=1024, overlap=0.35)
    results = []
    for (patch, x_offset, y_offset) in patches:
        dets = model(patch)            # {class, confidence, bbox}
        for det in dets:
            if det.confidence > conf_threshold:
                # Map bbox back to full image coordinates
                det.bbox = offset_bbox(det.bbox, x_offset, y_offset)
                results.append(det)
    return annotate_image(image, results)
\end{lstlisting}
